\textcolor{black}{\section{Conclusiones}}

\noindent Para el Dy$^{3+}$, se compararon los espectros obtenidos en dos matrices
vítreas distintas. En ambos se observaron bandas de emisión
características, aunque sus posiciones e intensidades relativas
presentaron variaciones. Esto sugiere que la matriz anfitriona influye
en el perfil de emisión del disprosio.

Para el Sm$^{3+}$, se compararon las mediciones del mismo vidrio de
fosfato de zinc realizadas con dos aparatos diferentes. Los máximos de
emisión aparecieron en longitudes de onda cercanas entre ambas
mediciones, por lo que las principales bandas se conservaron al cambiar
el instrumento. Las diferencias observadas pueden relacionarse con la
resolución, calibración y respuesta de cada aparato; por ello, las
intensidades absolutas no se compararon directamente.

En ambos casos, las bandas observadas se asociaron con transiciones
electrónicas características de los iones de tierras raras. Las
incertidumbres de longitud de onda se estimaron a partir de la anchura a
media altura de las bandas y las de número de onda se obtuvieron por
propagación. Estas incertidumbres son estimaciones espectrales y no
representan la incertidumbre instrumental.

Para el Tb$^{3+}$ se encontró que todos nuestros cálculos, considerando el rango de incertidumbre, de los niveles de transición que caracterizan su fotoluminiscencia están dentro de lo reportado en la literatura \cite{Tb-literatura}. A pesar de no tener un equipo altamente especializado, la configuración manual entre el láser y la fibra çoptica nos dio excelentes resultados.

Para el Eu$^{3+}$ se encontró casi la misma situación, solo hay una discrepancia con el decaimiento del $^5D_1$ al $^7F_6$, pues en la literatura \cite{Eu_Pr-literatura} se reporta que, para esa longitud de onda (aproximadamente), es del nivel $^5D_0$ al $^7F_4$.

Para el Pr$^{3+}$ también solo hay una discrepancia con un decaimiento en específico. Por la proximidad de estas dos intensidades y por el rango de incertidumbres, se encontró un mismo nivel de transcición: del $^3P_0$ al $^3F_2$. Se encontró \cite{Eu_Pr-literatura} que uno de ellos debería de caer al nivel $^3H_6$, aunque claro, con una longitud de onda un poco distinta.